% SPDX-License-Identifier: GPL-3.0-or-later
% Copyright (C) 2026 Alessandro Greco
\chapter{Глоссарий терминов Tor}\label[appendice]{cap:A-glossario}

Этот глоссарий определяет технические термины, используемые во всём русском издании. Для каждой записи также указана глава, в которой соответствующее понятие вводится.

{\small
\begin{longtable}{@{}>{\raggedright\arraybackslash}p{4.1cm}>{\raggedright\arraybackslash}p{8.0cm}r@{}}
\toprule
Термин & Значение & Гл. \\
\midrule
\endhead
\bottomrule
\endfoot
анализ трафика (traffic analysis) & Получение информации из временных характеристик, объёма и адресации коммуникаций без чтения их содержимого. & \ref{cap:01-anonimato} \\
анонимность (anonymity) & Состояние, при котором субъект остаётся неидентифицируемым внутри множества субъектов, называемого множеством анонимности. & \ref{cap:01-anonimato} \\
Sybil-атака & Представление сети множества идентичностей, например многочисленных ретрансляторов, чтобы увеличить вероятность выбора атакующего. & \ref{cap:26-attacchi} \\
центр каталогов (directory authority) & Один из небольшого числа серверов, голосующих за консенсус и подписывающих его. & \ref{cap:12-directory} \\
мост (bridge) & Ретранслятор, отсутствующий в публичном консенсусе и используемый как точка входа пользователями, сталкивающимися с цензурой. & \ref{cap:16-censura} \\
канал (channel) & TLS-соединение между двумя соседними узлами Tor, по которому передаются ячейки нескольких цепей. & \ref{cap:07-canali} \\
ячейка (cell) & Единица передачи в канале: обычно 514 байт, а для некоторых команд — переменной длины. & \ref{cap:08-celle} \\
ослеплённый ключ (blinded key) & Ключ, выводимый из идентификационного ключа onion-сервиса для заданного периода и скрывающий адрес сервиса от HSDir. & \ref{cap:14-onion-services} \\
ключ семейства (family key) & Ключ, используемый ретрансляторами одного оператора для подтверждения принадлежности к одному семейству. & \ref{cap:22-relay} \\
Circuit ID & Идентификатор цепи, локальный для одного канала; меняется на каждом переходе. & \ref{cap:08-celle} \\
цепь (circuit) & Маршрут через несколько ретрансляторов с послойным шифрованием, строящийся по одному переходу за раз. & \ref{cap:09-circuiti} \\
консенсус (consensus) & Подписанный центрами каталогов документ со списком ретрансляторов, их флагами и весами. & \ref{cap:12-directory} \\
управление перегрузкой (congestion control) & Регулирование объёма данных цепи в полёте в зависимости от измеренного времени прохождения туда и обратно. & \ref{cap:11-stream} \\
управление потоком (flow control) & Предотвращение перегрузки получателя отправителем с помощью окон и сообщений SENDME. & \ref{cap:11-stream} \\
корреляция трафика (traffic correlation) & Распознавание того, что два наблюдаемых потока в разных точках сети принадлежат одной коммуникации. & \ref{cap:26-attacchi} \\
дескриптор (descriptor) & Документ, посредством которого ретранслятор или onion-сервис публикует ключи, адреса и свойства. & \ref{cap:12-directory} \\
выходной ретранслятор & Последний ретранслятор цепи, открывающий соединение к назначению. & \ref{cap:06-architettura} \\
семейство ретрансляторов (family) & Группа ретрансляторов под общим управлением, которые не должны одновременно присутствовать в одной цепи. & \ref{cap:13-percorso} \\
флаг (flag) & Свойство, присваиваемое ретранслятору directory authorities, например \texttt{Guard}, \texttt{Exit} или \texttt{Stable}. & \ref{cap:12-directory} \\
прямая секретность (forward secrecy) & Свойство, при котором компрометация долгосрочного ключа не позволяет расшифровать ранее записанный трафик. & \ref{cap:05-crittografia} \\
guard, entry guard & Первый ретранслятор клиентских цепей, выбираемый на длительный срок и многократно используемый. & \ref{cap:13-percorso} \\
handshake & Обмен сообщениями, в ходе которого две стороны согласуют ключи, например ntor или ntor-v3. & \ref{cap:09-circuiti} \\
hidden service directory (HSDir) & Ретранслятор, хранящий и предоставляющий дескрипторы onion-сервисов. & \ref{cap:14-onion-services} \\
промежуточный ретранслятор (middle relay) & Ретранслятор в середине цепи, который не знает одновременно исходного клиента и назначения. & \ref{cap:06-architettura} \\
изоляция потоков (stream isolation) & Назначение разных потоков разным цепям, например на основе учётных данных SOCKS. & \ref{cap:18-stem} \\
протокол связи (link protocol) & Версионированный протокол канала, согласуемый ячейкой \cella{VERSIONS}. & \ref{cap:07-canali} \\
ретрансляционное сообщение (relay message) & Содержимое ячейки \cella{RELAY}, адресованное определённому переходу, например \texttt{BEGIN} или \texttt{DATA}. & \ref{cap:08-celle} \\
microdescriptor & Компактный дескриптор, содержащий только сведения о ретрансляторе, необходимые клиенту. & \ref{cap:12-directory} \\
модель угроз (threat model) & Явное описание противников и угроз, которым система должна противостоять. & \ref{cap:04-modello-minaccia} \\
несвязываемость (unlinkability) & Невозможность для наблюдателя определить, принадлежат ли два действия одному субъекту. & \ref{cap:01-anonimato} \\
onion-сервис (ранее hidden service) & Сервис, доступный через Tor, чей адрес кодирует его открытый ключ; ранее назывался hidden service. & \ref{cap:14-onion-services} \\
padding & Заполняющие данные в ячейках или канале, используемые для маскировки размеров и временных характеристик. & \ref{cap:07-canali} \\
веса пропускной способности (bandwidth weights) & Коэффициенты консенсуса, балансирующие выбор ретрансляторов для разных позиций цепи. & \ref{cap:13-percorso} \\
подключаемый транспорт (pluggable transport, PT) & Отдельная программа, преобразующая Tor-трафик для противодействия распознаванию или блокировке, например obfs4, Snowflake или WebTunnel. & \ref{cap:16-censura} \\
политика выхода (exit policy) & Упорядоченный набор правил, описывающий назначения, которые принимает выходной ретранслятор. & \ref{cap:22-relay} \\
доказательство выполнения работы (PoW) & Вычислительная задача, которую клиент решает, чтобы получить приоритет при доступе к onion-сервису во время атаки. & \ref{cap:15-dos} \\
предложение (proposal) & Нумерованный документ, используемый для обсуждения и проектирования изменений протокола Tor. & \ref{cap:31-futuro} \\
псевдоним (pseudonym) & Идентификатор, отличный от реального имени человека и позволяющий связывать действия без раскрытия личности. & \ref{cap:01-anonimato} \\
точка представления (introduction point) & Ретранслятор, выбранный onion-сервисом для получения начальных запросов на установление контакта. & \ref{cap:14-onion-services} \\
точка встречи (rendezvous point) & Ретранслятор, выбранный клиентом, в котором соединяются цепи клиента и onion-сервиса. & \ref{cap:14-onion-services} \\
ретранслятор (relay, onion router) & Узел сети Tor, пересылающий ячейки. & \ref{cap:06-architettura} \\
SENDME & Сообщение, подтверждающее получение данных цепи и повторно открывающее окно управления потоком. & \ref{cap:11-stream} \\
подпротокол (subprotocol) & Версионированное семейство функций протокола, объявляемое ретрансляторами и клиентами, например \texttt{Link} или \texttt{Relay}. & \ref{cap:30-arti} \\
поток (stream) & Поток прикладных данных, например TCP-соединение, передаваемый внутри цепи. & \ref{cap:11-stream} \\
vanguard (защитный промежуточный ретранслятор) & Ретранслятор промежуточного уровня, фиксируемый на определённый период для дополнительной защиты сторожевого узла onion-сервиса. & \ref{cap:13-percorso} \\
распознавание сайтов по трафику (website fingerprinting) & Распознавание посещённой страницы только по размерам и временным характеристикам зашифрованного трафика. & \ref{cap:26-attacchi} \\
\end{longtable}
}
